%% safe_numbers_seip.tex
%%
%% Copyright 2026 Matt Borland
%% Distributed under the Boost Software License, Version 1.0.
%% https://www.boost.org/LICENSE_1_0.txt
%%
%% ICSE Software Engineering in Practice (SEIP) submission draft.
%% Built on the IEEEtran conference skeleton (bare_conf.tex, M. Shell).
%% Compile with: latexmk -pdf safe_numbers_seip.tex
%%
%% Every figure in this paper is traceable to a file in the repository:
%%   RQ1  test/cve_corpus/analysis/{corpus.csv,results_by_cwe.csv}
%%   RQ2  doc/modules/ROOT/data/benchmarks-*/{unsigned,signed,float}.json
%%        (Table IV is regenerated by doc/paper/host_summary.py)
%%   RQ3  doc/modules/ROOT/pages/benchmarks.adoc (device sections)
%% Do not edit a number here without re-deriving it from those sources.

\documentclass[10pt,conference]{IEEEtran}
\IEEEoverridecommandlockouts % allow \thanks in conference mode for the corresponding author note

\usepackage{cite}
\usepackage{amsmath,amssymb}
\usepackage{booktabs}
\usepackage{array}
\usepackage{graphicx}
\usepackage{listings}

\makeatletter
% IEEEtran's \@makecaption breaks multi-line captions placed above a listing
\long\def\lst@makecaption#1#2{%
    \sbox\@tempboxa{\normalfont\footnotesize #1.\nobreakspace\nobreakspace #2}%
    \ifdim\wd\@tempboxa>\hsize
    {\normalfont\footnotesize\noindent #1.\nobreakspace\nobreakspace #2\par}%
    \else
        \ifCLASSOPTIONconference
        \hbox to\hsize{\hfil\box\@tempboxa\hfil}%
    \else
        \hbox to\hsize{\box\@tempboxa\hfil}%
    \fi
    \fi
    \vskip\belowcaptionskip}
\makeatother

\usepackage{xcolor}
\usepackage[hyphens]{url}
\usepackage{balance}
\usepackage{comment}

\lstset{
    language=C++,
    basicstyle=\ttfamily\scriptsize,
    keywordstyle=\bfseries,
    commentstyle=\itshape\color{black!60},
    columns=fullflexible,
    keepspaces=true,
    breaklines=true,
    frame=single,
    framerule=0.3pt,
    xleftmargin=4pt,
    xrightmargin=4pt,
    aboveskip=4pt,
    belowskip=4pt,
    abovecaptionskip=2pt,
    belowcaptionskip=6pt,
    showstringspaces=false
}

\newcommand{\mynote}[3]
{{\color{#1} \fbox{\bfseries\sffamily\scriptsize#2}
{\small\textsf{\emph{#3}}}}~}
\newcommand{\todo}[1]{\mynote{red}{TODO}{#1}}

\newcommand{\lib}{SafeNumbers}
\newcommand{\cmp}{Boost.SafeNumerics}
\newcommand{\code}[1]{\texttt{#1}}

\hyphenation{Safe-Numbers Safe-Numerics}

\begin{document}

    \title{\lib{}: Design and Empirical Evaluation of Checked Fixed-Width Numeric Types for C++,\\
    from Disclosed CVEs to Heterogeneous Accelerators}

    \author{\IEEEauthorblockN{Matt Borland\thanks{Corresponding author: Matt Borland (borlanm@purdue.edu).}}
    \IEEEauthorblockA{Purdue University, West Lafayette, IN, USA\\
    The C++ Alliance\\
    ORCID: 0009-0005-8183-4254\\
    Email: borlanm@purdue.edu}
    \and
    \IEEEauthorblockN{Jenna DiVincenzo}
    \IEEEauthorblockA{Purdue University, West Lafayette, IN, USA\\
    ORCID: 0000-0003-3029-2617\\
    Email: jennad@purdue.edu}
    }

    \maketitle

    \begin{abstract}

        Integer overflow, underflow, truncating conversion, and division by zero yield values
        that look valid but are wrong, and they remain common vulnerabilities in C and C++ software.
        Sanitizers typically run only in test builds, and trapping compiler flags can respond to a
        fault only by aborting.
        Overflow compiler intrinsics are non-portable, and must be written by hand at every call site.
        The checked-integer library in Boost follows C++ promotion rules by default, so a checked
        8-bit addition can evaluate to 256, and no other existing C++ libraries run inside GPU kernels.
        We present \lib{}, a C++20 library of checked fixed-width integer, floating-point, and range-bounded types.
        It rejects implicit and mixed-width conversions at compile time, and every operation is
        \code{constexpr}, so a fault during constant evaluation fails the build.
        Its central design decision decouples the detection of a fault, implemented once per
        operation and proven sound in Why3, from the response to it.
        This decoupling lets the same types run in GPU kernels using CUDA or SYCL, where a device fault is
        recorded and response occurs on the host.
        We evaluate the library against 140 Common Vulnerabilities and Exposures (CVE) entries
        from the National Vulnerability Database (NVD), covering four integer weakness classes and
        a negative control.
        It prevents 126 of the 130 in-scope faults (96.9\%, 95\% Wilson interval 92.4\% to 98.8\%),
        including every entry in the four primary classes.
        Across seven host targets, the median per-operation overhead is 1.00x to 1.18x of native
        arithmetic, and on a vectorizing SYCL device a branch-free deferred response brings
        compute-bound floating-point chains from up to 235x performance loss of native to
        1.1x to 2.9x while still reporting every error.

    \end{abstract}

    \begin{IEEEkeywords}
        integer overflow; software security; C++; safe arithmetic; empirical study;
        benchmarking; GPU computing; formal verification
    \end{IEEEkeywords}

    \IEEEpeerreviewmaketitle

    \section{Introduction}
    \label{sec:intro}

    Arithmetic errors, which include overflow past the largest value of a type, underflow below its smallest value, lossy conversions between types that truncate values, and division by zero, are pervasive even in mature, well-tested C++ code bases \cite{cwetop25, dietz2012} because such operations are valid and yield ordinary values rather than errors in the language. Since these values are functionally incorrect, they lead to miscalculations, corrupted program logic, or, at worst, critical software vulnerabilities. For example, a computation on a buffer size may cause the size to wrap around to a small number and result in a heap overflow or an out-of-bounds access when that value sizes an allocation or indexes an array.
%
% C++ makes these faults easy to write and hard to find, because its arithmetic rules turn
% them into ordinary values rather than errors.
% Unsigned arithmetic is specified to wrap, so an overflowed allocation size is well-defined
% behavior that no undefined-behavior tool will report.
% Signed overflow is undefined behavior, so an optimizer may assume that it never happens,
% including inside a check written after the fact to detect it.
% Integer promotion and the usual arithmetic conversions silently change operand types
% before an operation: narrow unsigned operands are widened to a signed \code{int}, and a
% signed operand combined with an unsigned one of the same width is converted to unsigned.
% The comparison \code{-1 < 1U} is therefore false, because \code{-1} is converted to the
% largest \code{unsigned int} before the comparison takes place.
%
    In practice, C++ software developers attempt to detect arithmetic errors in their code with sanitizers, compiler built-ins/intrinsics, or checked-integer libraries. But, each of these solutions have significant limitations when applied to modern C++ software.

    \emph{Sanitizers and their trapping flags check only some builds and behaviors, or respond only by aborting.}
    The primary sanitizer for C++, UndefinedBehaviorSanitizer (UBSan)~\cite{ubsan} leaves unsigned-overflow checks as opt-in making it easy to miss detection of such issues, such as in the case of the OpenSSH challenge-response handler (see Section \ref{sec:background} for more details). Furthermore, UBSan works by instrumenting a build to report faults as they occur at run time and is typically only used during testing as a result. Therefore, any gap in test coverage can let arithmetic errors slip into production undetected. Even if a developer configures a sanitizer to trap in production, the only response it supports on detection is to abort---turning a vulnerability into a denial of service with no path to recovery.

    \emph{Overflow intrinsics are precise but manual.}
    The GCC and Clang overflow intrinsics (\code{\_\_builtin\_*\_overflow}, e.g. \code{\_\_builtin\_add\_overflow(a, b, \&result)})~\cite{gccmanual} are built-in functions that perform an arithmetic operation and return true when an overflow occurs (false otherwise). They compile to the processor's carry and overflow flags, so space and performance costs are low and result precision is high. Additionally, the returned boolean result allows developers to respond to overflows however they like in code. Unfortunately, all of these benefits come at the cost of developer effort and portability: the check and error handling must be manually written at every operation that could fault (hopefully none are missed!) and the intrinsics are compiler extensions rather than portable C++.

    \emph{No checked-integer library reaches floating point types, 128-bit types, or device code.}
    Checked-integer libraries, such as \cmp{}~\cite{ramey2019}, SafeInt~\cite{safeint}, and Intel's safe-arithmetic library~\cite{intelsafe}, move the check into a type, so that no operation on a checked value goes unchecked automatically.
%\cmp{}~\cite{ramey2019}'s default policy operands follow C++ promotion before the check. So, for example,
% \code{safe<uint8\_t>(255) + 1} evaluates to 256 as an \code{int} and the 8-bit overflow is not reported where it occurs.
%Its range types check construction but not arithmetic results, it has no floating-point
%or 128-bit types, and its default response to a runtime fault is to throw, which device
%code cannot do.
    Unfortunately, no C++ library that we are aware of, offers checked floating-point types, 128-bit types, or device code.%, or executes inside a GPU kernel, so floating-point code, 128-bit code, and kernel code are left without a checked type.

    \emph{C++ software now runs on GPUs.} Modern C++ code increasingly runs on GPUs and other accelerators. Such device code, through say CUDA or SYCL, compiles into kernels that execute across thousands of parallel threads and is consequently constrained in ways host code is not. In particular, exceptions are unavailable because a checking scheme that throws them does not compile in a kernel. Additionally, performance depends on parallel threads executing the same instructions in compiler vectorizing loops. Any scheme that branches out of them early on arithmetic error can significantly degrade performance. Thus, mitigations that wish to cover whole C++ code bases must also work inside device code and cannot adopt the aforementioned traditional checking schemes.

    Taken together, these gaps define what a practical mitigation must do.
    It must
    (i)~protect production builds at every operation for integer, floating-point, and 128-bit types, not only in a tested subset or at the call sites a programmer remembered or for basic integer types;
    (ii)~check each operation at the width the programmer declared;
    (iii)~reject unsafe conversions at compile time;
    (iv)~let the program choose how to respond to a fault; and
    (v)~run in device code.
    No existing approach meets all five.

    This paper presents \lib{}~\cite{safenumbers},\footnote{\lib{} was developed with support
    from The C++ Alliance and is intended for proposal to Boost. It is not an official Boost
    library at the time of writing, although its namespace, \code{boost::safe\_numbers},
        anticipates that step.}
    a header-only C++20 library of checked numeric types designed to meet all five
    requirements.
    It provides fixed-width integers (\code{u8} through \code{u128} and \code{i8} through
    \code{i128}), floating-point types (\code{f32} and \code{f64}), and range-bounded types,
    and each of its design decisions answers one requirement.
    The check belongs to the type rather than to a build configuration or a call site, so it
    is compiled into every build and applied to every operation on every input~(i); adopting
    the library therefore changes declarations and the conversions at type boundaries, not
    every arithmetic expression.
    Every operation returns the type of its operands, so it is checked at the width the
    programmer declared, and the bounded types check every arithmetic result against their
    range, not only construction~(ii).
    Implicit and mixed-width conversions are rejected at compile time, and because every
    operation is \code{constexpr}, a fault during constant evaluation also fails to compile~(iii).
    \lib{}'s central design decision separates the \emph{detection} of a fault, implemented once per
    operation, from the \emph{response} to it, which can throw, saturate, return an optional
    or a flagged value, terminate, or invoke a user-defined handler, chosen per call site or
    per type~(iv).
    Because detection does not commit to a response, the same types run unchanged in CUDA
    and SYCL kernels, where a device fault is recorded and rethrown on the host, and a
    branch-free response lets a loop vectorize~(v).
%Together, these decisions move the point at which an arithmetic fault is caught: into the shipped program, on whatever input an attacker supplies, and, for conversion faults, into the compiler, before the program ships at all.

    We evaluate \lib{} both formally and empirically. A checked type is only as trustworthy as its detection logic,
%, and an edge case such as \code{INT\_MIN / -1} is easy to miss; since no test suite can enumerate every pair of 128-bit operands,
    so we verified the soundness of our detection algorithms in the Why3 deductive verification platform~\cite{why3}. We also assess \lib{}'s relevance, cost, and reach through answering the following research questions:
%Three claims must hold in practice for that to matter, and each research question tests one: the faults the types detect must be the ones behind real vulnerabilities, the checks must be cheap enough that developers keep them in release builds rather than reverting to the test-only position of a sanitizer, and the responses that the decoupling enables must run inside device kernels at an acceptable cost.

    \begin{itemize}
        \item \textbf{RQ1 (Relevance).} To what extent do checked fixed-width and domain-bounded
        types prevent the integer-arithmetic weaknesses behind real, disclosed Common Vulnerabilities and Exposures entries (CVEs)?
        \item \textbf{RQ2 (Cost).} What runtime overhead does always-on checking impose, relative
        to native arithmetic and to \cmp{}?
        \item \textbf{RQ3 (Reach).} Which error-response mechanisms let the same checked types
        execute in exception-free, massively parallel device code, and at what cost?
    \end{itemize}

    To answer RQ1, we sourced 140 CVEs from the National Vulnerability Database (NVD) across four arithmetic weakness classes and a
    negative control, reduced each CVE to its root-cause
    computation, and wrote that computation once with built-in types and once with \lib{}. The library prevents 126 of the 130 in-scope faults, including every entry in the four
    primary classes.
    To answer RQ2, we benchmarked every operation at every width on seven operating-system
    and architecture targets against native arithmetic, \cmp{}, SafeInt, and Intel's
    safe-arithmetic library. We found a median overhead of at most 1.16x of native
    arithmetic through 64 bits on 64-bit hosts.
    To answer RQ3, we timed memory-bound and compute-bound kernels under CUDA on an NVIDIA
    GPU and under SYCL on a CPU device, and found that most checks add no measurable cost to
    memory-bound GPU kernels, while on the SYCL device a branch-free response that records
    faults and reports them afterward keeps compute-bound floating-point chains within 3x of
    native arithmetic, down from up to 235x.

    In summary our contributions, all open source under the Boost Software License~\cite{safenumbers},
    are:
    \begin{enumerate}
        \item \lib{}, a header-only C++20 library of checked 8- to 128-bit integer,
        floating-point, and bounded types whose detection of a fault is decoupled from the
        response to it, so one set of types serves throwing host code, exception-free code,
        vectorized loops, and CUDA and SYCL kernels, with the detection algorithms proven sound in
        195 Why3 lemmas (Section~\ref{sec:library}).
        \item A reproducible corpus of 140 NVD-sourced CVEs across four integer weakness classes
        and a negative control, each reduced to a self-contained test with machine-readable
        metadata, and a script that regenerates every corpus statistic
        (Section~\ref{sec:design-rq1}).
        \item A per-operation cost model for checked arithmetic when no fault occurs, across
        seven host targets and two heterogeneous back ends, with a semantic probe matrix and a
        cost comparison against three other checked-integer libraries
        (Sections~\ref{sec:results-rq2} and~\ref{sec:results-rq3}).
        \item An ablation over the library's own response policies showing that in vectorizable
        loops the cost is set by the kind of response, one that yields a value or one that exits
        the loop early, rather than by the detection code, which makes the decoupling the enabling
        mechanism for vectorized loops and device kernels (Sections~\ref{sec:results-rq2}
        and~\ref{sec:discussion}).
        \item Lessons for practitioners on where checked arithmetic could be used and how to
        choose a response (Section~\ref{sec:discussion}).
    \end{enumerate}

    The remainder of this paper is organized as follows.
    Section~\ref{sec:background} describes the integer weakness classes and the C++ rules
    that make them easy to write.
    Section~\ref{sec:library} presents the design of \lib{}.
    Section~\ref{sec:design} describes the study design for each research question, and
    Section~\ref{sec:results} reports the results.
    Section~\ref{sec:discussion} discusses lessons for practitioners,
    Section~\ref{sec:threats} the threats to validity, and Section~\ref{sec:related} related
    work.
    Section~\ref{sec:conclusion} concludes and outlines future work.

%% OLD INTRO
    \begin{comment}
        The recent push toward memory-safe software, including guidance from national
        cybersecurity agencies to migrate away from memory-unsafe languages~\cite{cisa2024},
        has focused attention on pointers, lifetimes, and bounds.
        Many memory-safety failures, however, begin earlier, in arithmetic that produces a value
        which looks valid but is wrong.
        A buffer size that wraps around to a small number, or a length that underflows past zero
        to a large one, raises no error where it is computed.
        The fault surfaces only later, when that value sizes an allocation or indexes an array,
        as a heap overflow or an out-of-bounds access.

        MITRE's Common Weakness Enumeration (CWE), a catalog of software weakness types, groups
        these arithmetic faults into four classes (Section~\ref{sec:background}):
        integer overflow or wraparound (CWE-190), where a result exceeds the largest value of its
        type;
        integer underflow (CWE-191), where it falls below the smallest;
        incorrect conversion between numeric types (CWE-681), such as a signed count compared
        against an unsigned bound or a 64-bit length truncated to 32 bits;
        and division by zero (CWE-369).
        These weaknesses are common in practice.
        Integer overflow has appeared in the CWE Top 25 Most Dangerous Software Weaknesses for
        years~\cite{cwetop25}, and Dietz et al.\ showed more than a decade ago that it is
        pervasive even in mature, well-tested C and C++ code bases~\cite{dietz2012}.
        All four classes also recur among Common Vulnerabilities and Exposures (CVE) entries, the
        public identifiers assigned to disclosed vulnerabilities, which the National Vulnerability
        Database (NVD) annotates with their CWE class.

        C++ makes these faults easy to write and hard to find, because its arithmetic rules turn
        them into ordinary values rather than errors.
        Unsigned arithmetic is specified to wrap, so an overflowed allocation size is well-defined
        behavior that no undefined-behavior tool will report.
        Signed overflow is undefined behavior, so an optimizer may assume that it never happens,
        including inside a check written after the fact to detect it.
        Integer promotion and the usual arithmetic conversions silently change operand types
        before an operation: narrow unsigned operands are widened to a signed \code{int}, and a
        signed operand combined with an unsigned one of the same width is converted to unsigned.
        The comparison \code{-1 < 1U} is therefore false, because \code{-1} is converted to the
        largest \code{unsigned int} before the comparison takes place.

        These rules apply wherever C++ runs, and C++ increasingly runs on GPUs and other
        accelerators.
        Through CUDA and SYCL, the same C++ source compiles into kernels that execute across
        thousands of parallel lanes.
        Device code imposes constraints that host code does not: exceptions are unavailable, and
        performance depends on lanes executing the same instructions and on the compiler
        vectorizing loops, which a data-dependent early exit can prevent.
        A checking scheme that throws does not compile in a kernel, and one that branches out of
        a hot loop can cost more than the computation it protects.
        Any mitigation intended to cover a whole C++ program must therefore also work inside
        device code.

        \emph{Sanitizers and trapping flags check only some builds, or respond only by aborting.}
        UndefinedBehaviorSanitizer (UBSan)~\cite{ubsan} instruments a build to report faults as
        they occur.
        It is typically used as a testing configuration, so it finds a fault only when a test
        input reaches it, whereas the input behind a CVE is chosen by an attacker after the
        software ships.
        Its unsigned-overflow check is opt-in, because unsigned wrapping is defined behavior.
        The \code{-ftrapv} flag~\cite{gccmanual} can remain enabled in release builds, but it
        covers signed overflow only, while many allocation-size faults, including the OpenSSH
        example in Section~\ref{sec:background}, are unsigned.
        Where \code{-ftrapv} is deployed, or a sanitizer is configured to trap in production, the
        only response is to abort the process, which turns a memory-corruption vulnerability into
        a denial of service and leaves no path to recovery.

        \emph{Overflow intrinsics are precise but manual.}
        The GCC and Clang overflow intrinsics (\code{\_\_builtin\_*\_overflow})~\cite{gccmanual}
        compile to the processor's carry and overflow flags and return a flag that the caller can
        act on however it chooses.
        That flexibility has a price: the check must be written at every operation that could
        fault, so one forgotten call site is the vulnerability, and the intrinsics are compiler
        extensions rather than portable C++.

        \emph{\cmp{} checks the promoted type rather than the declared one.}
        Checked-integer libraries move the check into a type, so that no operation on a checked
        value goes unchecked.
        \cmp{}~\cite{ramey2019}, the checked-integer library in Boost, takes this approach, but
        under its default policy operands follow C++ promotion before the check, so
        \code{safe<uint8\_t>(255) + 1} evaluates to 256 as an \code{int} and the 8-bit overflow is
        not reported where it occurs.
        Its range types check construction but not arithmetic results, it has no floating-point
        or 128-bit types, and its default response to a runtime fault is to throw, which device
        code cannot do.

        \emph{No checked-integer library reaches floating point, 128-bit types, or device code.}
        SafeInt~\cite{safeint} and Intel's safe-arithmetic library~\cite{intelsafe} also
        provide checked integer types for host code.
        Neither, nor any other C++ library we are aware of, offers checked floating-point or
        128-bit types, or executes inside a GPU kernel, so floating-point code, 128-bit code, and
        kernel code are left without a checked type.

        Taken together, these gaps define what a practical mitigation must do.
        It must
        (i)~be present in shipped builds and at every operation, not only in tests or at the call
        sites a programmer remembered;
        (ii)~check each operation at the width the programmer declared;
        (iii)~reject unsafe conversions at compile time;
        (iv)~let the program choose how to respond to a fault; and
        (v)~run in device code.
        No existing approach meets all five.

        This paper presents \lib{}~\cite{safenumbers},\footnote{\lib{} was developed with support
        from The C++ Alliance and is intended for proposal to Boost. It is not an official Boost
        library at the time of writing, although its namespace, \code{boost::safe\_numbers},
            anticipates that step.}
        a header-only C++20 library of checked numeric types designed to meet all five
        requirements.
        It provides fixed-width integers (\code{u8} through \code{u128} and \code{i8} through
        \code{i128}), floating-point types (\code{f32} and \code{f64}), and range-bounded types,
        and each of its design decisions answers one requirement.
        The check belongs to the type rather than to a build configuration or a call site, so it
        is compiled into every build and applied to every operation on every input~(i); adopting
        the library therefore changes declarations and the conversions at type boundaries, not
        every arithmetic expression.
        Every operation returns the type of its operands, so it is checked at the width the
        programmer declared, and the bounded types check every arithmetic result against their
        range, not only construction~(ii).
        Implicit and mixed-width conversions are rejected at compile time, and because every
        operation is \code{constexpr}, a fault during constant evaluation also fails the
        build~(iii).
        The central decision separates the \emph{detection} of a fault, implemented once per
        operation, from the \emph{response} to it, which can throw, saturate, return an optional
        or a flagged value, terminate, or invoke a user-defined handler, chosen per call site or
        per type~(iv).
        Because detection does not commit to a response, the same types run unchanged in CUDA
        and SYCL kernels, where a device fault is recorded and rethrown on the host, and a
        branch-free response lets a hot loop vectorize~(v).
        A checked type is only as trustworthy as its detection logic, and an edge case such as
        \code{INT\_MIN / -1} is easy to miss; since no test suite can enumerate every pair of
        128-bit operands, we prove the detection algorithms sound in the Why3 deductive
        verification platform~\cite{why3}.

        Together, these decisions move the point at which an arithmetic fault is caught: into the
        shipped program, on whatever input an attacker supplies, and, for conversion faults, into
        the compiler, before the program ships at all.
        Three claims must hold in practice for that to matter, and each research question tests
        one: the faults the types detect must be the ones behind real vulnerabilities, the checks
        must be cheap enough that developers keep them in release builds rather than reverting to
        the test-only position of a sanitizer, and the responses that the decoupling enables must
        run inside device kernels at an acceptable cost.

        \begin{itemize}
            \item \textbf{RQ1 (Relevance).} To what extent do checked fixed-width and domain-bounded
            types prevent the integer-arithmetic weaknesses behind real, disclosed CVEs?
            \item \textbf{RQ2 (Cost).} What runtime overhead does always-on checking impose, relative
            to native arithmetic and to \cmp{}?
            \item \textbf{RQ3 (Reach).} Which error-response mechanisms let the same checked types
            execute in exception-free, massively parallel device code, and at what cost?
        \end{itemize}

        To answer RQ1, we sourced 140 CVEs from the NVD across the four weakness classes and a
        negative control (CWE-682, Incorrect Calculation), reduced each to its root-cause
        computation, and wrote that computation once with built-in types and once with \lib{};
        the library prevents 126 of the 130 in-scope faults, including every entry in the four
        primary classes.
        To answer RQ2, we benchmarked every operation at every width on seven operating-system
        and architecture targets against native arithmetic, \cmp{}, SafeInt, and Intel's
        safe-arithmetic library, and found a median overhead of at most 1.16x of native
        arithmetic through 64 bits on 64-bit hosts.
        To answer RQ3, we timed memory-bound and compute-bound kernels under CUDA on an NVIDIA
        GPU and under SYCL on a CPU device, and found that most checks add no measurable cost to
        memory-bound GPU kernels, while on the SYCL device a branch-free response that records
        faults and reports them afterward keeps compute-bound floating-point chains within 3x of
        native arithmetic, down from up to 235x.

        Our contributions, all open source under the Boost Software License~\cite{safenumbers},
        are:
        \begin{enumerate}
            \item \lib{}, a header-only C++20 library of checked 8- to 128-bit integer,
            floating-point, and bounded types whose detection of a fault is decoupled from the
            response to it, so one set of types serves throwing host code, exception-free code,
            vectorized loops, and CUDA and SYCL kernels, with the detection algorithms proven sound in
            195 Why3 lemmas (Section~\ref{sec:library}).
            \item A reproducible corpus of 140 NVD-sourced CVEs across four integer weakness classes
            and a negative control, each reduced to a self-contained test with machine-readable
            metadata, and a script that regenerates every corpus statistic
            (Section~\ref{sec:design-rq1}).
            \item A per-operation cost model for checked arithmetic when no fault occurs, across
            seven host targets and two heterogeneous back ends, with a semantic probe matrix and a
            cost comparison against three other checked-integer libraries
            (Sections~\ref{sec:results-rq2} and~\ref{sec:results-rq3}).
            \item An ablation over the library's own response policies showing that in vectorizable
            loops the cost is set by the kind of response, one that yields a value or one that exits
            the loop early, rather than by the detection code, which makes the decoupling the enabling
            mechanism for vectorized loops and device kernels (Sections~\ref{sec:results-rq2}
            and~\ref{sec:discussion}).
            \item Lessons for practitioners on where checked arithmetic could be used and how to
            choose a response (Section~\ref{sec:discussion}).
        \end{enumerate}S

        The remainder of this paper is organized as follows.
        Section~\ref{sec:background} describes the integer weakness classes and the C++ rules
        that make them easy to write.
        Section~\ref{sec:library} presents the design of \lib{}.
        Section~\ref{sec:design} describes the study design for each research question, and
        Section~\ref{sec:results} reports the results.
        Section~\ref{sec:discussion} discusses lessons for practitioners,
        Section~\ref{sec:threats} the threats to validity, and Section~\ref{sec:related} related
        work.
        Section~\ref{sec:conclusion} concludes and outlines future work.
    \end{comment}

    \section{Background and Motivation}
    \label{sec:background}

    \subsection{Integer Weaknesses as Root Causes}\label{subsec:integer-weaknesses-as-root-causes}

    We consider four weakness classes from the NVD whose root cause is a fixed-width integer computation.
    \emph{CWE-190, Integer Overflow or Wraparound}, leads to the classic allocation-size fault.
    In CVE-2002-0639 the OpenSSH challenge-response handler computed
    \code{nresp * sizeof(char*)} in 32 bits, an attacker-controlled \code{nresp} of $2^{30}$
    wrapped the product to zero, and the subsequent writes overflowed the heap.
    \emph{CWE-191, Integer Underflow}, goes the other way, typically a length minus a header
    size that wraps from zero to a huge value.
    \emph{CWE-681, Incorrect Conversion Between Numeric Types}, covers signedness confusion
    and truncation: a negative signed length compared against an unsigned bound passes the
    check because the comparison converts it to a large positive value.
    \emph{CWE-369, Divide by Zero}, occurs where an untrusted divisor reaches a division or modulo.
    In each, a representable but wrong value propagates to a memory operation, and the CVE is
    filed against that operation; the arithmetic upstream is the root cause. The negative control class used is \emph{CWE-682, Incorrect Calculation}.

    \subsection{Why the Language Does Not Help}

    Consider the comparison \code{-1 < 1U}.
    Intuitively the result is \code{true}.
    In C++ the signed operand converts to \code{unsigned int}, so the expression is
    \code{false}. The conversion is well defined, and it is diagnosed only when \code{-Wsign-compare} is enabled.
    The hardware already computes carry and overflow flags on every arithmetic instruction, what is missing is a type that consults them.

    \section{\lib{}}
    \label{sec:library}

    This section presents the design of \lib{}: each decision, the gap in the existing
    mitigations it closes (Section~\ref{sec:intro}), and what it costs.
    The full reference is available online~\cite{safenumbers}.

    \subsection{Explicit Over Implicit}

    The library prohibits implicit conversion in every direction: between safe and built-in
    types, between safe types of different widths, and between signed and unsigned safe types.
    These are all compile-time errors, as are construction from \code{bool} and unary minus
    on unsigned types.
    This closes the promotion gap of \cmp{}.
    An operation on two \code{u8} values is checked
    as an 8-bit operation and yields a \code{u8}.
    A signedness confusion (CWE-681) no
    longer compiles, so it is found even if there is no test input that triggers this issue.
    The cost is verbosity at type boundaries, where a \code{static\_cast} is required, but
    every conversion becomes visible in the source.
    These type requirements also enables simpler verification.
    Every operation stays within a single ring $\mathbb{Z}/2^n\mathbb{Z}$, which removes widening and narrowing from the proof obligations.

    \subsection{Fail as Early as Possible}

    Every operation is \code{constexpr}.
    The default error path contains a \code{throw} expression, which is not permitted in a
    constant expression, so an overflow during constant evaluation becomes a compile-time
    error rather than a run-time exception.
    The developer writes one expression and receives the strongest guarantee the context
    allows: a build failure where the operands are known, or a check where they are not.

    \subsection{Detection Decoupled from Response}
    \label{sec:library-policy}

    Handling an arithmetic fault has two parts, detecting it and responding to it.
    The mitigations in Section~\ref{sec:intro} either fix the response or leave it to each
    call site, so none fits both a server that must recover and a kernel that cannot throw.
    \lib{} implements detection once per operation, classifying the outcome into an
    \code{error\_kind} (\code{overflow}, \code{underflow}, \code{divide\_by\_zero}, and for
    floating point \code{nan\_operation} and \code{invalid\_operation}), and routes the
    classification to a response chosen independently.
    The detection code, and the proofs that cover it, are therefore shared by every response.
    Three API surfaces expose that choice to the developer.

    \emph{Default operators} throw a typed standard exception
    (\code{std::overflow\_error}, \code{std::underflow\_error}, \code{std::domain\_error}).
    The throw itself is contained in a function annotated as cold, and forced non-inline so that an operator's body
    is only the arithmetic, the check, and a branch.

    \emph{Call-site families}, named after Rust's primitive-type methods~\cite{rustnum},
    select the response per operation without changing the type:
    \code{checked\_add} returns \code{std::optional}, \code{saturating\_add} clamps,
    \code{overflowing\_add} returns a \code{std::pair} of the (possibly wrapped) value and a flag,
    \code{strict\_add} immediately terminates with \code{std::exit},
    and \code{widening\_add} promotes to the next width (e.g.\ \code{u32 + u32} yields \code{u64}).

    \emph{Type-level policies} carry the response in the type for every operator:
    \code{sat\_u8} saturates, \code{strict\_i32} terminates, and
    \code{basic\_u32<Handler>} invokes a user-defined stateless handler whose
    \code{on\_error(error\_kind, value, message)} return value becomes the operation's
    result.
    Only responses whose result is the operand type may live in the type, so a value of a
    safe type is always exactly its numeric value.
    The \code{noexcept} specification of every operator follows the handler, so a
    non-throwing handler yields non-throwing arithmetic.

    \subsection{Bounded Domains as an Additional Precondition}

    The bounded types express a value's legal domain in the type
    (\code{bounded\_uint<Min, Max>}, \code{bounded\_int}, and \code{bounded\_float}) and
    enforce it on construction and on every arithmetic result, unlike the range types of \cmp{}, which check only construction.
    This is modeled on Ada's range subtypes~\cite{spark}, and
    it reaches faults that would be legal within a type width.
    A correct addition whose result is one past the end of a 150-element table does not
    overflow, but \code{bounded\_uint<0, 149>} rejects it when the value is formed,
    before a bad-memory access could occur.
    Whenever a bounded value is converted out of its type, the declared domain is still provided
    to the optimizer as an assumption, so checks on the receiving type that the domain makes
    impossible are deleted (Section~\ref{sec:results-rq2}).

    \subsection{Device Execution}
    \label{sec:library-device}

    All integer, floating-point, and bounded types compile with NVCC and with Intel's
    oneAPI DPC++ compiler for execution inside CUDA and SYCL kernels, so host and device code
    share one set of checked types.
    Device code cannot throw, so only the response half of handling a fault changes.
    A \code{device\_error\_context} contains a shared error record (CUDA \code{\_\_managed\_\_}
    memory or a SYCL \code{device\_global}) into which the first faulting work item writes the
    file, line, thread identifier, expression, and exception type.
    A subsequent call to \code{synchronize()} on the host waits for the kernel,
    and throws the same standard exception as if the operation had been executed on the host.
    For hot kernels a second response exists, the \code{overflowing\_*} families produce a
    value and a flag with no branch, and a \code{deferred\_errors} accumulator collects the
    flags across a chain and raises them once per work item (Listing~\ref{lst:deferred}),
    trading the precision of the report for a straight-line kernel body.

    \begin{lstlisting}[float=t,caption={A SYCL kernel using the deferred response. Detection is
    unchanged; the response is one branch per work item.},label={lst:deferred}]
boost::safe_numbers::device_error_context ctx {q};

q.parallel_for(sycl::range<1>(n), [=](sycl::id<1> idx)
{
    const auto i {idx[0]};
    boost::safe_numbers::deferred_errors errs;

    // unwrap() folds each flag in without branching
    auto x {errs.unwrap(overflowing_add(in[i], in[i]))};
    x = errs.unwrap(overflowing_mul(x, in[i]));

    errs.report(__FILE__, __LINE__); // one branch
    out[i] = x;
});

// rethrows on the host if any item reported
ctx.synchronize();
    \end{lstlisting}

    \subsection{Soundness of the Detection Logic}
    \label{sec:library-verification}

    Following the verification of GMP's low-level routines in Why3~\cite{rieuhelft2017,
        melquiond2018}, the overflow, underflow, and edge-case detection for every fundamental
    operation was re-expressed in WhyML and proven with Why3~\cite{why3}.
    The unsigned theory clones Why3's generic bit-vector module so one proof covers every width from 8 to 128 bits.
    The signed theory proves the two's-complement sign-bit identities, and the asymmetric
    cases \code{INT\_MIN / -1} and \code{-INT\_MIN}.
    Lastly, the floating-point theory proves that the IEEE 754 result classifier assigns the correct category to every input.
    In total, we verify 195 lemmas.
    The C++ source is a manual translation of the verified algorithms, so the claim is
    soundness of the detection logic, not verification of the shipped source.

    \section{Study Design}
    \label{sec:design}

    \subsection{RQ1: CVE Corpus}
    \label{sec:design-rq1}

    \subsubsection{Sourcing}
    CVEs were fetched from the NVD via their REST API 2.0~\cite{nvdapi}, by CWE identifier, on 19 August 2026.
    Four primary categories (CWE-190, CWE-191, CWE-681, CWE-369) were sampled with thirty relevant examples each.
    CWE-682, Incorrect Calculation, serves as a negative control with twenty examined entries.
    Many of its defects are in-range logic errors that no arithmetic check can catch, so we
    expect a lower rate.

    \subsubsection{Reproduction}
    For each CVE we extracted the computation that the disclosure and its patch identify as
    the root cause and wrote a self-contained test set with three sections: a native part that computes
    the expression with built-in types and asserts the specific silent wrong value the
    disclosure describes, a safe part that computes the identical expression with the library
    types and asserts the fault (a thrown exception of the expected type, or a compile
    failure), and a bounded part that declares the value's documented domain with a bounded
    type and asserts rejection at construction.
    Where the original width was platform dependent we model the width the advisory
    publishes; 96 of the 130 in-scope entries model 32 bits, and the 30 CWE-681 entries mix
    signed and unsigned operands.

    \subsubsection{Outcome categories}
    Each examined CVE receives one outcome: prevented at run time (the safe computation throws
    at the faulting operation), prevented at compile time (the bug is ill-formed using the library and is
    committed as a compile-fail test), prevented by a bound (the declared domain rejects the
    input before the arithmetic), not prevented (a genuine arithmetic CVE whose faulting value
    is in range, counted in the denominator), or out of scope (not a fixed-width integer fault
    on reduction, recorded with a reason and reported separately).
    Each test carries a machine-readable metadata block, and a dependency-free script
    cross-checks the blocks against the test outcomes and regenerates every corpus table,
    exiting non-zero on any inconsistency, so the corpus is validated in continuous integration.

    \subsubsection{Statistics}
    The primary metric per category is the prevention rate, prevented over in scope, with a
    95\% Wilson score interval~\cite{wilson1927}, reported once with the default checked types alone and once including the bounded types.
    The native baseline detects zero faults as expected, so every prevented case is a
    discordant pair and the exact McNemar test~\cite{mcnemar1947} reduces to
    $p = \min(1, 2 \cdot 0.5^{b})$ for $b$ prevented.
    We reproduce the root-cause computation, not the original program, so the claim is that
    the fault would have surfaced at the faulting operation, not that the whole program would
    have run correctly (Section~\ref{sec:threats}).

    \subsection{RQ2: Host Benchmark Grid}
    \label{sec:design-rq2}

    The benchmarks use Google Benchmark~\cite{gbench}.
    Each benchmark fills a vector of one million random values once, then for every operation
    evaluates \code{op(vec[i], vec[i+1])} across the vector, passing each result through
    \code{DoNotOptimize} to ensure the computation is not elided.
    Code is compiled with \code{-O2}, or MSVC equivalent.
    Integer operands are drawn from a uniform distribution over $[1, w-1]$ for a $w$-bit type and
    sorted in descending order so that all five operations stay on the non-overflowing path.
    Floating-point operands are drawn uniformly from $[1, 100]$.
    The measurement therefore reflects the cost of the checks that production code pays on
    every operation, not the cost of throwing, which is paid only on a fault.
    We run five repetitions and report the median time per element.

    Three contenders are compared at each width: the built-in type as baseline, the library
    type, and \cmp{}'s \code{safe<T>} with its default template arguments.
    \cmp{} has no floating-point and no 128-bit type, so those tables only have two contenders.
    The grid ran on seven GitHub Actions runners: Linux x64 and ARM64 (GCC 14), Linux x86
    32-bit (GCC 14, \code{-m32}), macOS ARM64 (Apple Clang), and Windows x64, ARM64, and x86
    32-bit (MSVC), all at C++20.

    \subsubsection{Semantics of the comparators}
    RQ2 is a question of cost-effectiveness rather than cost alone, because \cmp{} checks
    promoted rather than declared widths by default, so we first establish what each
    comparator detects with ten semantic probes, each a single arithmetic situation with
    runtime operands (Table~\ref{tab:probes}), expressed idiomatically in \lib{}, \cmp{},
    Microsoft SafeInt~\cite{safeint}, and Intel safe-arithmetic~\cite{intelsafe}, all in their
    default configuration.
    A driver compiles every (library, probe) pair, records a compile failure as a compile-time
    rejection with its diagnostic, and otherwise runs the binary and records its output line.
    Intel safe-arithmetic tracks intervals in the type and proves containment at compile time,
    with a runtime check only where a raw value enters an interval type through
    \code{safe::match}, so its probes admit the operands through \code{match}.

    \subsubsection{Cross-library cost}
    A second benchmark runs the built-in type, \lib{}, \cmp{}, SafeInt, and Intel
    safe-arithmetic on the same operands at 32 and 64 bits.
    The Intel library admits both operands through \code{match} and performs the proven operation unchecked (unsigned
    subtraction on equal intervals is not expressible and is therefor omitted).

    \subsubsection{Ablation over response policies}
    A third benchmark isolates the response half of the architecture, testing the claim that
    the choice of response, not the shared detection code, sets the cost.
    For 8-, 32-, 64-, and 128-bit unsigned and 32- and 128-bit signed widths it runs the five
    operations on the built-in type and on seven responses: the default throwing operators,
    the saturating and strict aliases, two user-defined handlers (a \code{noexcept} handler
    returning the wrapped fallback, and one throwing a user exception type), the deferred
    \code{overflowing\_*} family with its flag folded without a branch, and the
    \code{checked\_*} family unwrapped with a fallback.
    A further pair of rows measures what a declared domain provides.
    The default operators on operands validated into a bounded type at ingress, against a built-in loop over the same
    one-byte operands.

    \subsubsection{Two kernel shapes}
    The host grid times a \emph{scalar} shape in which every element's result passes through
    \code{DoNotOptimize}, which forces a round trip per element and prevents vectorization for
    every contender alike.
    The cross-library and ablation benchmarks additionally time a \emph{stream} shape,
    \code{out[i] = a[i] op b[i]} into an output array with one memory clobber after the pass,
    the shape real element-wise loops have, which lets the optimizer vectorize whichever
    contenders it can.
    Both ran locally on an Apple M4 Max (macOS 26.6, Apple Clang 21.0, \code{-O2}, five
    repetitions, median reported), with inputs on the non-violating path as in the grid.


    \subsection{RQ3: Device Benchmarks}
    \label{sec:design-rq3}

    Two dedicated benchmarks, one per back end, time identical kernels for each safe type and
    its basis type, so the ratio isolates the cost of the checks.
    Inputs are chosen so no check fires, and a \code{device\_error\_context} covering the run
    aborts the benchmark if one does.
    Two kernel shapes are measured: \emph{streaming}, \code{out[i] = a[i] op b[i]} over
    4\,194\,304 elements, the shape of real element-wise kernels, memory bound on a discrete
    GPU; and a \emph{compute-bound chain} in which 65\,536 threads each evolve one
    register-resident value through 256 cycles of eight to nine checked operations.
    The integer chain mixes in division and modulo so that the optimizer cannot fold the
    native baseline's wrapping additions; floating-point chains are per operation, since IEEE
    arithmetic is not associative.
    Every family also runs the non-throwing forms against the same baseline: a saturating
    addition streaming row, saturating and overflowing integer chains, and deferred
    floating-point chains whose flags are accumulated and folded into the result.
    Every timed region is verified afterwards: safe and native results must match bit for
    bit.

    The CUDA benchmark ran on an NVIDIA GeForce RTX 3060 (sm\_86, 28 SMs) with CUDA 13.1 on
    Ubuntu 24.04 at \code{-O2}.
    The SYCL benchmark ran on the OpenCL CPU device of an Intel Core i9-11900K with oneAPI
    icpx 2026.0 at \code{-O3}.
    Because icpx defaults to \code{-ffp-model=fast}, which grants the optimizer floating-point
    associativity and would let the native chains fold, the build pins
    \code{-ffp-model=precise} so both contenders run IEEE arithmetic exactly as under NVCC.

    \section{Results}
    \label{sec:results}

    \subsection{RQ1: Prevention of Disclosed CVEs}
    \label{sec:results-rq1}

    Table~\ref{tab:rates} gives the disposition and prevention rate of every category.
    Of 140 examined CVEs, 10 were out of scope, all in the CWE-682 control, for reasons
    recorded in the corpus.
    The 130 in-scope entries span 2002 to 2022 across products from the Linux kernel to
    LibTIFF, QEMU, and FreeType.
    The exact McNemar $p$ is $1.9\times10^{-9}$ for each primary category, $0.031$ for the
    control, and $2.4\times10^{-38}$ pooled.

    \begin{table}[t]
        \centering
        \caption{Disposition and prevention rate by category. Nothing examined is dropped;
        out-of-scope entries are reported and removed from the denominator. ``Default'' counts the
        checked types alone; ``Bnd.'' adds cases prevented once the domain is declared with a
        bounded type. The rate, in percent, is on the latter with its 95\% Wilson score interval.}
        \label{tab:rates}
        \footnotesize
        \setlength{\tabcolsep}{3pt}
        \begin{tabular}{@{}lrrrrrl@{}}
            \toprule
            Category & Exam. & Out & In & Default & Bnd. & Rate (95\% CI) \\
            \midrule
            CWE-190 & 30 & 0 & 30 & 29 & 30 & 100.0 (88.6 to 100.0) \\
            CWE-191 & 30 & 0 & 30 & 30 & 30 & 100.0 (88.6 to 100.0) \\
            CWE-681 & 30 & 0 & 30 & 30 & 30 & 100.0 (88.6 to 100.0) \\
            CWE-369 & 30 & 0 & 30 & 30 & 30 & 100.0 (88.6 to 100.0) \\
            CWE-682 (ctrl.) & 20 & 10 & 10 & 3 & 6 & 60.0 (31.3 to 83.2) \\
            \midrule
            All & 140 & 10 & 130 & 122 & 126 & 96.9 (92.4 to 98.8) \\
            \bottomrule
        \end{tabular}
    \end{table}

    \subsubsection{Primary categories}
    All four primary categories are prevented at 100\% including the bounded arm, with a
    Wilson lower bound of 88.6\% at $n = 30$.
    The mechanism that fires differs by category.
    CWE-190, CWE-191, and CWE-369 are prevented at runtime by the throwing operators.
    In CWE-681, 17 of 30 entries are prevented at \emph{compile time}: the root cause is a
    comparison or arithmetic between a signed and an unsigned quantity, and porting the
    expression to \code{i32} and \code{u32} makes it ill-formed; the remaining 13 are
    truncating conversions that throw on the explicit narrowing cast.
    One CWE-190 entry is prevented only through the bounded arm: in CVE-2018-13785 (libpng)
    an oversized but representable image width overflows a row-factor calculation, and
    declaring libpng's documented width limit as a \code{bounded\_uint} rejects it at the
    boundary, before the downstream faults the runtime checks would also have caught.

    \subsubsection{The negative control}
    CWE-682 behaves as a control should.
    The default types prevent three of ten in-scope cases (divide-by-zero and shift faults
    the NVD tagged as calculation errors), and declared domains prevent three more
    (CVE-2011-3062, CVE-2011-1573, CVE-2018-11790), each a wrong index, length, or offset used
    against a buffer of independently known size.
    The remaining four stay unprevented because their wrong result is a valid in-domain value:
    three are integer divisions that truncate where the formula needed rounding or a ceiling
    (CVE-2016-7433, CVE-2017-8932, CVE-2018-14439), and one is an index one past the end with
    no independently known bound to declare (CVE-2018-20999).
    No bound applies, which is exactly the library's limit and why the control is informative.

    \subsubsection{Answer to RQ1}
    Among disclosed CVEs whose root cause reduces to a fixed-width integer arithmetic fault,
    the checked types convert the fault into a typed, catchable failure in every sampled case
    of the four primary classes, bounded domains extend prevention to in-range faults the
    operation-level checks cannot see, and the pooled rate is 96.9\% (92.4\% to 98.8\%),
    bounded below by the honest failures in the control.
    The explicit-conversion rule catches 17 of the conversion faults with no test input at all.

    \subsection{RQ2: Host Overhead}
    \label{sec:results-rq2}

    Table~\ref{tab:host} summarizes the grid as the geometric mean, median, and maximum of
    the per-operation ratio across each type class; the full 58-row tables per platform are
    in the documentation.

    \begin{table*}[t]
        \centering
        \caption{Host overhead relative to the built-in type of the same width, non-violating
        path, median of five repetitions. Each cell is geometric mean / median / maximum of the
        per-operation ratio across the class (40 operations for int 8 to 64, 10 for 128, 8 for
        float). \cmp{} has neither 128-bit nor floating-point types.}
        \label{tab:host}
        \footnotesize
        \setlength{\tabcolsep}{5pt}
        \begin{tabular}{@{}llcccc@{}}
            \toprule
            Platform & Compiler & SafeNumbers int 8 to 64 & SafeNumerics int 8 to 64 & SafeNumbers 128 & SafeNumbers float \\
            \midrule
            Linux x64        & GCC 14      & 1.14 / 1.16 / 1.66 & 1.25 / 1.25 / 4.72  & 1.10 / 1.10 / 1.81 & 1.18 / 1.20 / 1.29 \\
            Linux ARM64      & GCC 14      & 1.12 / 1.02 / 1.71 & 1.18 / 1.06 / 10.60 & 1.01 / 1.03 / 1.04 & 1.47 / 1.46 / 1.88 \\
            Linux x86 32-bit & GCC 14      & 1.34 / 1.18 / 4.32 & 1.29 / 1.10 / 4.56  & 1.48 / 1.24 / 3.46 & 1.06 / 1.05 / 1.34 \\
            macOS ARM64      & Apple Clang & 1.07 / 1.08 / 1.37 & 0.95 / 0.96 / 1.04  & 1.16 / 1.07 / 3.05 & 1.01 / 1.01 / 1.04 \\
            Windows x64      & MSVC        & 1.08 / 1.03 / 1.50 & 1.04 / 1.01 / 1.50  & 1.32 / 1.22 / 4.99 & 1.02 / 1.02 / 1.09 \\
            Windows ARM64    & MSVC        & 1.03 / 1.00 / 2.36 & 1.07 / 1.00 / 2.99  & 1.61 / 1.19 / 7.73 & 1.51 / 1.44 / 1.86 \\
            Windows x86 32-bit & MSVC      & 1.26 / 1.16 / 5.69 & 1.08 / 1.03 / 2.00  & 1.97 / 2.15 / 4.09 & 2.66 / 3.83 / 4.73 \\
            \bottomrule
        \end{tabular}
    \end{table*}

    \subsubsection{Through 64 bits}
    On every 64-bit host the median per-operation overhead of the checked integer types lies
    between 1.00x and 1.16x, and the geometric mean between 1.03x and 1.14x.
    Division and modulo are at parity almost everywhere (on Linux x64 all 16 division and
    modulo rows through 64 bits are at 1.00x), because the only check is a compare against a
    divisor already in flight.
    Addition and subtraction cost about 1.2x on GCC x64 and are at parity on MSVC.
    Multiplication is the most expensive check, up to 1.7x with GCC and 2.36x for signed
    64-bit multiplication on Windows ARM64, because the overflow test needs a wide product or
    a division; that one cell is the only miss of the design target of less than 2x on
    64-bit hosts.

    \subsubsection{32-bit targets}
    The two 32-bit targets are where the cost is visible: on Linux x86 32-bit six of the 40
    integer rows exceed 2x, all at 8 and 16 bits (3.0x to 4.3x), where GCC's code for the
    narrow overflow intrinsics is poor; on Windows x86 32-bit, 64-bit multiplication reaches
    5.69x because product and check are both emulated in 32-bit registers, and the
    floating-point classifier costs 1.7x to 4.7x.

    \subsubsection{128-bit and floating-point}
    The 128-bit types cost 1.01x to 1.48x by geometric mean on GCC and Clang hosts; their
    multiplication check, computed from 64-bit limbs because there is no wider arithmetic to
    lift into, is the one genuinely expensive check in the library, peaking at 3.05x to 7.73x
    on every host except Linux x64 and ARM64, where GCC keeps it at parity.
    The floating-point classifiers cost 1.18x on GCC x64, are at parity on Apple Clang and
    MSVC x64, and reach 1.5x under GCC on Linux ARM64 and MSVC on Windows ARM64.

    \subsubsection{Against the comparators: semantics}
    Table~\ref{tab:probes} gives the outcome of the ten semantic probes; the four libraries
    embody three models.
    \lib{} and SafeInt check at the declared width at runtime and agree on every runtime fault
    (P1 to P5 and P7); at the type boundary \lib{} rejects mixed signedness, mixed width, and
    \code{bool} construction at compile time, while SafeInt compares mixed signedness correctly
    (P6), accepts \code{bool}, and rejects mixed widths only through an overload ambiguity.
    \cmp{} under its default promotion widens the 8-bit overflow to 256 (P1), lets a bounded
    result escape its range (P10), and throws on a mixed-signedness comparison instead of
    comparing.
    Intel safe-arithmetic rejects at compile time everything its interval logic can see and is
    the one comparator whose bounded type enforces on results (P10), but it accepts
    \code{bool}, and its division interval ignores the sign of the quotient, so
    \code{INT\_MIN / -1} compiles and returns \code{INT\_MIN} (P4).
    Only \lib{} converts every one of the ten situations into a compile-time or runtime fault.

    \begin{table*}[t]
        \centering
        \caption{Semantic probes with runtime operands, default configuration of each library.
        CT: rejected at compile time. RT: runtime fault (exception). Other cells give the value the
        program printed. Intel s.-a. is Intel safe-arithmetic.}
        \label{tab:probes}
        \footnotesize
        \setlength{\tabcolsep}{6pt}
        \begin{tabular}{@{}lllll@{}}
            \toprule
            Probe & \lib{} & \cmp{} & SafeInt & Intel s.-a. \\
            \midrule
            P1 \code{u8} $255 + 1$ & RT & widened (256) & RT & CT \\
            P2 \code{u32} $0 - 1$ & RT & RT & RT & CT \\
            P3 \code{u32} $2^{31} \cdot 2$ & RT & RT & RT & CT \\
            P4 \code{i32} \code{INT\_MIN} $/$ $-1$ & RT & RT & RT & wrong (\code{INT\_MIN}) \\
            P5 \code{u32} $x / 0$ & RT & RT & RT & CT \\
            P6 \code{i32}$(-1) <$ \code{u32}$(1)$ & CT & RT & correct (true) & CT \\
            P7 \code{u32} 300 to \code{u8} & RT & RT & RT & CT \\
            P8 \code{u8} $+$ \code{u32} & CT & widened (2) & CT (ambiguous) & widened (2) \\
            P9 construct from \code{bool} & CT & accepted (1) & accepted (1) & accepted (1) \\
            P10 $[0, 100]$: $60 + 50$ & RT & escaped (110) & no bounded type & CT \\
            \bottomrule
        \end{tabular}
    \end{table*}

    \subsubsection{Against the comparators: cost}
    Table~\ref{tab:libs} places the four libraries on the same operands and stream shape;
    cost follows the response class rather than the library.
    The three libraries that detect at the declared width and throw, \lib{}, \cmp{}, and
    SafeInt, form one class: on unsigned 32-bit addition, subtraction, and multiplication
    they cost 2.6x to 4.4x, on unsigned 64-bit 1.3x to 1.6x (with one \cmp{} addition
    outlier at 2.5x), and all three are at parity on unsigned division and modulo.
    Intel safe-arithmetic is the cheapest contender at 0.95x to 1.34x on 32 bits, but its
    runtime work is not smaller, only differently placed: two range checks on the inputs with
    a fallback value, rather than one check on the result with an exit.
    Two controls isolate the effect.
    Giving \code{match} a throwing fallback moves Intel to 3.6x to 5.1x, the class of the
    throwing libraries, and removing its runtime check gives 0.95x to 1.3x, so the checks
    themselves cost at most 0.3x, about what our value-returning \code{checked} family costs
    on addition and subtraction (Table~\ref{tab:ablation}).
    Interval reasoning is genuinely cheaper only for multiplication, where a bound check on two
    inputs vectorizes and a check on the high half of the product does not (1.34x against 4.0x
    to 5.2x for our unsigned saturating, overflowing, and checked forms).
    The contracts also differ: Intel enforces a declared input domain and cannot express a
    full-width addition whose overflow is detected, so the two coincide only when the domain is
    known, which is what this benchmark constructs.

    Signed 32-bit arithmetic is where the comparison improved the library.
    In our first run \lib{} cost 10.1x to 13.3x on signed addition, subtraction, and
    multiplication, against 4.5x for SafeInt, because Clang declined to inline every signed
    throwing operator (cost 320 against a threshold of 225) once the exception, its message,
    and its source location were constructed inline.
    Outlining that error path into the cold, non-inlined function of
    Section~\ref{sec:library-policy} restored inlining for all twenty signed operators;
    Table~\ref{tab:libs} reports the library after the fix, 4.0x to 6.1x, the class of
    SafeInt (3.4x to 4.5x), while \cmp{} remains at 10.0x to 13.1x.

    \begin{table}[t]
        \centering
        \caption{Cross-library cost at 32 bits, ratio to the built-in type, stream shape, Apple M4
        Max with Apple Clang 21 at \code{-O2}, median of five repetitions. Intel safe-arithmetic
        performs two boundary range checks and an unchecked proven operation.}
        \label{tab:libs}
        \footnotesize
        \setlength{\tabcolsep}{3pt}
        \begin{tabular}{@{}llcccccc@{}}
            \toprule
            Type & Library & Add & Sub & Mul & Div & Mod & Gmean \\
            \midrule
            \code{u32} & SafeNumbers & 2.60 & 3.63 & 4.20 & 0.98 & 1.01 & 2.09 \\
            \code{u32} & Boost.SafeNumerics & 2.91 & 3.67 & 4.36 & 0.97 & 1.02 & 2.15 \\
            \code{u32} & Microsoft SafeInt & 2.68 & 3.68 & 4.07 & 1.04 & 1.01 & 2.11 \\
            \code{u32} & Intel safe-arithmetic & 0.97 & n/a & 1.34 & 1.05 & 1.05 & 1.09 \\
            \midrule
            \code{i32} & SafeNumbers & 3.99 & 5.13 & 6.05 & 0.98 & 1.02 & 2.62 \\
            \code{i32} & Boost.SafeNumerics & 9.99 & 13.13 & 10.61 & 1.81 & 0.99 & 4.78 \\
            \code{i32} & Microsoft SafeInt & 3.40 & 3.97 & 4.49 & 0.99 & 0.98 & 2.26 \\
            \code{i32} & Intel safe-arithmetic & 0.95 & 1.28 & 1.25 & 1.04 & 1.02 & 1.10 \\
            \bottomrule
        \end{tabular}
    \end{table}

    \subsubsection{Ablation over response policies}
    Table~\ref{tab:ablation} is the ablation on the stream shape, where the built-in 32-bit
    loops vectorize four lanes wide and the 64-bit loops two.
    The responses split into two classes, and the compiler's own diagnostics say why.
    Every response that must act at the exact faulting element (throw, a throwing handler, or
    \code{strict}'s termination) is an early exit and draws the remark ``cannot vectorize
    early exit loop with writes to memory''; these cost 2.8x to 5.1x for 32-bit addition,
    subtraction, and multiplication, falling to 1.3x to 1.4x at 64 bits where only two lanes
    are lost.
    Every response whose result is a value keeps the loop straight-line and vectorizes: the
    wrapping handler (0.74x to 1.32x), saturating addition and subtraction, which lower to the
    native saturating vector instructions (0.74x to 1.00x), and the deferred
    \code{overflowing} family and \code{checked} unwrapped with a fallback (0.96x to 1.29x for
    addition and subtraction).
    Multiplication is the exception: the unsigned overflow test needs the high half of the
    product, which has no 32-bit vector form on this target, so unsigned saturating,
    overflowing, and checked multiplication, and signed saturating multiplication, cost 4.0x
    to 5.2x, while the signed deferred and checked forms stay at 1.4x to 1.6x.
    Division and modulo are at parity for every response because the built-in loop does not
    vectorize them either.
    At 128 bits, where the basis type is a two-word class, the default operators and value
    responses stay within 1.25x (unsigned) and 1.6x (signed) except for multiplication, where
    the throwing and saturating responses cost 1.6x to 1.8x unsigned and 3.2x to 3.3x signed,
    and the deferred and checked forms 7x to 9x.

    The bounded ingress rows show what a declared domain buys.
    Because the bounded types hand their domain to the optimizer as an assumption whenever a
    value leaves the type, operands validated into \code{bounded\_uint<1,31>} and widened to
    \code{u32} carry the fact that they lie in $[1, 31]$, and Clang deletes every check that
    fact makes impossible before vectorizing: the default throwing operators cost 1.02x to
    1.36x on 32-bit addition and multiplication and 0.81x to 1.07x at 128 bits.
    Unsigned subtraction keeps its check (5.8x) because $[1, 31] - [1, 31]$ can underflow;
    the optimizer removes exactly what the declared facts justify.

    The result corrects an expectation we held going in: \code{strict}, modeled on Rust's
    panic on overflow, does not vectorize, because termination at the faulting element is an
    early exit exactly as a throw is (Rust's release-mode arithmetic vectorizes because it
    wraps, not because it panics).
    What vectorizes is detection whose response is a value, the property that makes the
    deferred family the enabling mechanism on devices (Section~\ref{sec:results-rq3}).
    Replacing a response with another of the same class leaves the cost where the class puts
    it: the user throwing handler matches the default operators for unsigned 32-bit (3.70x
    against 3.67x for addition), and the wrapping handler matches the built-in type.
    The signed user throwing handler (9.2x to 11.5x) is the inlining failure above again: its
    inline \code{throw} pushes the operator over Clang's threshold (235 against 225), so a user
    handler that throws should, like the library, call a cold, non-inlined function.
    In the scalar shape, where nothing vectorizes, all seven responses fall within 1.0x to
    1.2x of one another apart from cells affected by a harness artifact
    (Section~\ref{sec:threats}).

    \begin{table}[t]
        \centering
        \caption{Ablation over response policies, ratio to the built-in type, stream shape, same
        machine as Table~\ref{tab:libs}. Every row runs the same detection code; only the
        response differs. $^{\dagger}$Operands validated into \code{bounded\_uint<1,31>} or
        \code{bounded\_int<1,31>} at ingress and widened per element; ratio to a built-in loop
        over the same one-byte operands.}
        \label{tab:ablation}
        \footnotesize
        \setlength{\tabcolsep}{3pt}
        \begin{tabular}{@{}llcccccc@{}}
            \toprule
            Type & Response & Add & Sub & Mul & Div & Mod & Gmean \\
            \midrule
            \code{u32} & throwing (default) & 3.67 & 3.51 & 3.87 & 0.99 & 0.98 & 2.17 \\
            \code{u32} & saturating & 0.98 & 1.00 & 3.95 & 0.97 & 0.97 & 1.30 \\
            \code{u32} & strict & 3.54 & 3.31 & 3.74 & 0.98 & 0.99 & 2.12 \\
            \code{u32} & user handler, wrapping & 1.32 & 1.17 & 1.12 & 0.98 & 1.01 & 1.11 \\
            \code{u32} & user handler, throwing & 3.70 & 3.31 & 3.71 & 0.97 & 0.97 & 2.12 \\
            \code{u32} & overflowing (deferred) & 1.16 & 0.96 & 4.36 & 0.97 & 0.98 & 1.36 \\
            \code{u32} & checked (optional) & 1.29 & 1.21 & 5.16 & 0.98 & 1.00 & 1.51 \\
            \code{u32} & throwing, bounded ingress$^{\dagger}$ & 1.36 & 5.81 & 1.09 & 1.00 & 1.00 & 1.54 \\
            \midrule
            \code{i32} & throwing (default) & 3.85 & 3.92 & 4.69 & 1.03 & 1.05 & 2.38 \\
            \code{i32} & saturating & 1.00 & 0.74 & 4.19 & 1.02 & 1.08 & 1.28 \\
            \code{i32} & strict & 5.13 & 2.80 & 3.45 & 1.02 & 1.07 & 2.22 \\
            \code{i32} & user handler, wrapping & 0.87 & 0.74 & 0.74 & 1.06 & 1.04 & 0.88 \\
            \code{i32} & user handler, throwing & 11.53 & 10.32 & 9.20 & 1.02 & 1.04 & 4.10 \\
            \code{i32} & overflowing (deferred) & 1.00 & 0.96 & 1.57 & 1.05 & 1.08 & 1.11 \\
            \code{i32} & checked (optional) & 0.99 & 1.01 & 1.43 & 1.03 & 1.03 & 1.09 \\
            \code{i32} & throwing, bounded ingress$^{\dagger}$ & 1.24 & 0.90 & 1.02 & 0.97 & 0.98 & 1.02 \\
            \bottomrule
        \end{tabular}
    \end{table}

    \subsubsection{Answer to RQ2}
    On 64-bit hosts, always-on checking at the declared width costs a median of 1.00x to
    1.16x per operation through 64 bits, with division and modulo free and multiplication the
    most expensive; the exceptions (32-bit targets, 128-bit multiplication, and floating-point
    classification on some ARM64 compilers) are quantified above.
    Relative to \cmp{} the cost is comparable where the semantics are comparable, and \lib{}
    detects faults that library's default configuration silently absorbs.
    On a vectorizable loop the class of the response sets the cost regardless of library: a
    response that is a value vectorizes at 0.7x to 1.3x of native, and an early exit costs 3x
    to 6x at 32 bits.

    \subsection{RQ3: Device Execution}
    \label{sec:results-rq3}

    \subsubsection{Streaming}
    Table~\ref{tab:stream} gives selected streaming rows for both devices.
    On the GPU the checks are essentially free: the kernels are bound by memory bandwidth and
    the checks execute in the shadow of the loads; the rows above parity are the 8-bit types
    (up to 1.23x), \code{f64} division (1.44x), and 128-bit multiplication (6.8x).
    The SYCL CPU device tells the same story for 32- to 128-bit integers (within 11\%),
    but not for \code{f32} (1.3x to 1.7x) or below 32 bits (up to 1.9x at 16 bits):
    four million one-byte elements fit in cache, the kernels are compute bound, a CPU device
    executes a kernel by vectorizing it across work items, and the per-operation reporting
    branch inhibits exactly that, so the 8-bit rows reach 11.6x.
    \code{saturating\_add} compiles to the native saturating vector instructions, returning
    \code{u8} to parity and \code{i8} to 1.5x.

    \begin{table}[t]
        \centering
        \caption{Streaming kernels, ratio of safe time to basis-type time, selected rows.
        RTX 3060 under CUDA and Intel i9-11900K OpenCL CPU device under SYCL.}
        \label{tab:stream}
        \footnotesize
        \setlength{\tabcolsep}{4pt}
        \begin{tabular}{@{}lcccccc@{}}
            \toprule
            & \multicolumn{3}{c}{CUDA (RTX 3060)} & \multicolumn{3}{c}{SYCL (CPU device)} \\
            \cmidrule(lr){2-4}\cmidrule(lr){5-7}
            Type & Add & Mul & Sat.\ add & Add & Mul & Sat.\ add \\
            \midrule
            \code{u8}   & 1.05 & 1.05 & 1.02 & 5.47 & 11.60 & 0.97 \\
            \code{u32}  & 1.00 & 1.00 & 1.00 & 1.03 & 1.04  & 1.00 \\
            \code{u64}  & 1.00 & 1.00 & 1.00 & 1.03 & 1.04  & 1.00 \\
            \code{u128} & 1.00 & 6.77 & 1.00 & 1.00 & 1.05  & 0.99 \\
            \code{i8}   & 1.23 & 1.22 & 1.04 & 9.69 & 11.32 & 1.52 \\
            \code{i32}  & 1.00 & 1.00 & 1.00 & 1.00 & 1.00  & 0.96 \\
            \code{i64}  & 1.00 & 1.00 & 1.00 & 1.04 & 1.03  & 1.00 \\
            \code{f32}  & 1.00 & 1.00 & n/a  & 1.32 & 1.35  & n/a \\
            \code{f64}  & 1.00 & 1.00 & n/a  & 1.05 & 1.05  & n/a \\
            \bottomrule
        \end{tabular}
    \end{table}

    \subsubsection{Compute-bound integer chains}
    With the arithmetic pinned in registers the worst case becomes visible.
    On the GPU the throwing integer types cost 1.30x to 2.93x through 64 bits and 7.78x to
    10.30x at 128 bits, and the non-throwing forms buy little (1.21x to 2.96x) because a
    throwing operator already compiles to a well-predicted branch.
    On the CPU device the dominant cost is lost vectorization: where the baseline cannot
    vectorize either the throwing types cost 1.05x to 1.88x, and where it can (\code{i16},
    \code{i32}) the gap is the SIMD width times the check cost, 8.8x to 9.4x, with
    \code{u32} between the two at 4.3x.

    \subsubsection{Compute-bound floating-point chains}
    Table~\ref{tab:fchain} is the central result of RQ3.
    A floating-point chain's baseline operation is a single hardware instruction, so the
    fixed cost of the classifier multiplies it.
    On the GPU the throwing \code{f32} chains cost 2.2x (division) to 18.3x and the
    \code{f64} chains 1.5x to 4.0x, and the deferred family cuts \code{f32} to 2.7x to 5.4x
    with division at parity.
    On the CPU device the throwing chains lose vectorization entirely and cost 17x to 235x.
    The deferred family, whose \code{overflowing\_*} operations are branch-free, vectorizes
    with the baseline and lands the same chains at 1.05x to 2.85x, two orders of magnitude
    better, while still detecting every error at the boundary.

    \begin{table}[t]
        \centering
        \caption{Compute-bound floating-point chains, ratio to the basis type, throwing
        operators versus the deferred (overflowing) family.}
        \label{tab:fchain}
        \footnotesize
        \begin{tabular}{@{}llcccc@{}}
            \toprule
            & & \multicolumn{2}{c}{CUDA (RTX 3060)} & \multicolumn{2}{c}{SYCL (CPU device)} \\
            \cmidrule(lr){3-4}\cmidrule(lr){5-6}
            Type & Op & Throwing & Deferred & Throwing & Deferred \\
            \midrule
            \code{f32} & Add & 18.31 & 5.39 & 173.03 & 1.68 \\
            \code{f32} & Sub & 17.66 & 5.37 & 171.30 & 2.11 \\
            \code{f32} & Mul &  8.71 & 2.69 & 127.59 & 1.43 \\
            \code{f32} & Div &  2.22 & 1.05 &  50.01 & 1.13 \\
            \code{f64} & Add &  4.00 & 4.21 & 179.70 & 2.73 \\
            \code{f64} & Sub &  4.01 & 4.21 & 234.65 & 2.85 \\
            \code{f64} & Mul &  2.77 & 2.94 & 102.24 & 2.05 \\
            \code{f64} & Div &  1.49 & 1.55 &  17.00 & 1.05 \\
            \bottomrule
        \end{tabular}
    \end{table}

    \subsubsection{Answer to RQ3}
    Two responses built on the same detection code make device execution possible: a shared
    device-side error record with host-side rethrow, which preserves the throwing model across
    the kernel boundary, and a branch-free deferred family, which preserves vectorization.
    In the memory-bound regime typical of real kernels the GPU checks are free above 8 bits
    apart from 128-bit multiplication and \code{f64} division.
    In the compute-bound regime the throwing types cost 1.3x to 2.9x for integers through 64
    bits on a GPU, and the deferred family is what makes checked floating point viable on a
    vectorizing device, at 1.1x to 2.9x instead of up to 235x.

    \section{Discussion and Lessons for Practitioners}
    \label{sec:discussion}

    \subsection{Fixed Width Is the Point}
    The comparison with \cmp{} shows that per-operation cost is similar when the two libraries
    check the same thing.
    What differs is the semantics: promotion-based checking cannot see an 8-bit or 16-bit
    fault that the surrounding expression computes in \code{int}, and 96 of the 130 in-scope
    CVEs in our corpus are 32-bit faults that promotion would not have widened away on the
    platforms where they shipped, while the 30 signedness faults are ones promotion actively
    hides.
    SafeInt, the other fixed-width comparator, agrees with the library on every runtime probe.
    Practitioners choosing a checked type should ask at which width the check is applied,
    not only how fast it is.
    Within a width, division and modulo are free, addition and subtraction cost one predictable
    branch, and multiplication, which needs a wide product or a division, is the check to watch.

    \subsection{Decoupling Detection from Response Is the Enabling Mechanism}
    On the host, offering saturating, checked, and overflowing families beside the throwing
    operators looks like convenience; on a device it is the difference between a library that
    works and one that does not.
    Exceptions do not exist in a kernel, so the response must be relocated to the host
    (\code{device\_error\_context}) or deferred to a boundary (\code{deferred\_errors}), with
    the detection code untouched in both cases.
    The floating-point chains (Table~\ref{tab:fchain}) are the sharpest illustration: the same
    classifier, with the branch moved from every operation to once per work item, goes from
    unusable to under 3x on the vectorizing device, and the host ablation shows the same split
    in the vectorizer's diagnostics.
    We suggest the lesson applies to any checked-arithmetic design intended for
    performance-sensitive code: make the response a policy, and make at least one policy a
    value.

    \subsection{Bounded Domains Catch What Checks Cannot}
    The CWE-682 control doubles its prevented count when domains are declared.
    Any value later used as an index, offset, or length into a buffer of independently known
    size has a domain, and declaring it moves the failure from the memory operation to the
    point the value is formed; the bound must come from that independent quantity, never from
    the correct answer.

    \subsection{Costs of Adoption}
    The explicit-conversion rule adds casts at every boundary with legacy code, and the
    mixed-width prohibition forces a decision wherever C++ would have promoted.
    The corpus suggests that friction is the mechanism: 17 of the 30 CWE-681 entries were
    prevented because the port would not compile.

    \section{Threats to Validity}
    \label{sec:threats}

    \subsubsection{Construct validity}
    For RQ1, ``prevented'' means that the arithmetic root cause becomes a typed, catchable
    failure at the faulting operation.
    A thrown exception is strictly better than silent corruption, but is a denial of service
    if unhandled, and we do not claim the original program would have been immune.
    For RQ2 and RQ3, the ratio measures the non-violating path by construction; the cost of
    actually throwing is not measured and is deliberately out of scope.

    \subsubsection{Internal validity}
    Reproduction fidelity is the central threat to RQ1.
    We do not re-run the original
    programs but port only the operation the patch cites, at the width the advisory assumes,
    assert the specific disclosed wrong value, and commit every arm as a test that runs in CI.
    For RQ2, the host numbers come from shared CI runners and reflect a single run each.
    Five repetitions with median reporting stabilize the central figure but cannot remove
    the variance of a shared virtual machine, so differences of a few percent should be read
    as noise.
    The grid's scalar shape forces a per-element spill whose cost depends on stack layout.
    In one build the throwing \code{u32} and \code{u64} measured 4.3x to 4.9x because the AArch64
    load/store optimizer merged two spill reloads into one \code{ldp} straddling two prior
    stores, defeating store-to-load forwarding, disabling that pass returned them to 1.01x.
    An isolated outlier cell of Table~\ref{tab:host} may therefore reflect the harness rather
    than a check; the stream-shape comparisons do not share this exposure.
    The stream-shape ratios on vectorized 32-bit loops rest on baselines near 0.1 ns per
    element and vary by up to 30\% between runs, so their second digit is not meaningful.
    For RQ3, each device is a single machine and a single run.

    \subsubsection{External validity}
    The corpus is a purposive sample of CVEs whose root cause reduces to a self-contained
    snippet, so the prevention rates are conditional on that reducibility and are not rates
    over all CVEs tagged with a CWE; the out-of-scope funnel, the negative control, and the
    not-prevented cases carry the credibility, not a claim of random sampling.
    The device results come from one consumer GPU and one CPU OpenCL device.

    \subsubsection{Conclusion validity}
    Because $n$ is small and the observed rate is at the ceiling, the rate and its Wilson
    interval are the result; the McNemar $p$ values are ceilinged by design.

    \section{Related Work}
    \label{sec:related}

    \subsubsection{Understanding and finding integer faults}
    Dietz et al.~\cite{dietz2012} found integer overflow pervasive in mature C/C++, including
    intentional wrapping that complicates detection; their tool became UBSan's integer
    checks~\cite{ubsan}.
    KINT~\cite{wang2012}, RICH~\cite{brumley2007}, and SmartFuzz~\cite{molnar2009} find
    integer errors statically, by binary instrumentation, and by test generation; Coker and
    Hafiz~\cite{coker2013} and IntPatch~\cite{zhang2010} repair them by transformation.
    These tools find or repair faults in code already written; \lib{} keeps its checks on in
    the shipped program, and we evaluate it against the CVE record those studies characterized.

    \subsubsection{Compiler and language mechanisms}
    Rust~\cite{rustnum} panics on overflow in debug builds and wraps in release, with the
    \code{checked\_}, \code{saturating\_}, \code{wrapping\_}, and \code{overflowing\_} methods
    the library's call-site families follow by name; Ada and SPARK~\cite{spark} provide range
    subtypes, with GNATprove proving absence of overflow on Why3's back end.
    C++26 contracts~\cite{p2900} standardize the violation-handling model the library's
    response architecture anticipates; library hardening~\cite{p3471} and constrained
    numbers~\cite{p2993} are adjacent WG21 efforts.

    \subsubsection{Safe-integer libraries}
    \cmp{}~\cite{ramey2019} is the comparator throughout.
    Microsoft's SafeInt~\cite{safeint} is a widely deployed checked-integer template that
    throws or calls a handler.
    \code{bounded::integer}~\cite{bounded} tracks value bounds at compile time so that
    overflow is impossible by construction, at the cost of result types that widen.
    \code{type\_safe}~\cite{typesafe} rejects lossy and implicit conversions.
    Intel safe-arithmetic~\cite{intelsafe} carries an interval in the type, proves that every
    operation's result interval fits its destination at compile time, and admits runtime
    values only through a checked \code{match}; it is explicitly pre-release, and
    Section~\ref{sec:results-rq2} reports both what that model catches and one case it
    misses.
    None provides floating-point classification, 128-bit types, or device execution.

    \subsubsection{Verified arithmetic}
    The GMP verification in Why3~\cite{rieuhelft2017,melquiond2018} is the model for our
    proofs, Glesner~\cite{glesner2006} supplies the lifting argument for multiplication, and
    refinement types~\cite{vazou2014} inspired the bounded types' role as a precondition layer,
    though our bounds are checked dynamically.

    \section{Conclusion and Future Work}
    \label{sec:conclusion}

    We designed and implemented \lib{}, a C++20 library of checked fixed-width numeric types
    that keeps its checks in the shipped program, checks at the declared width, and decouples
    the detection of a fault from the response to it, and we evaluated each decision against a
    measurable claim.
    It prevents 126 of 130 in-scope CVEs; its median per-operation cost on 64-bit hosts is
    1.00x to 1.16x through 64 bits; and on CUDA and SYCL devices the same types run unchanged,
    with a branch-free deferred response bringing compute-bound floating-point chains on a
    vectorizing SYCL device from two orders of magnitude slower than native to under 3x.
    Building and evaluating the library together improved both: the evaluation exposed an
    inlining failure in the signed throwing operators, now fixed, and led the bounded types to
    hand their declared domain to the optimizer.
    Future work is a random corpus draw, a controlled-machine replication of the host grid
    with dispersion statistics, and a discrete-GPU SYCL run.

    \section*{Acknowledgment}
    The first author is an employee of The C++ Alliance, which funded the development of
    \lib{} and the evaluation reported here.
    The authors thank the C++ Alliance for its continued support of open-source C++ libraries and tools.
    Anthropic's Claude (Opus 5, Opus 5.5, Fable 5.1, and Sonnet 5.5 models) were used to draft and edit prose throughout this paper.
    The authors verified every claim and number against the artifacts cited in Section~\ref{sec:design}.

    \begin{thebibliography}{99}

        \bibitem{cisa2024}
        Cybersecurity and Infrastructure Security Agency and Federal Bureau of Investigation,
        ``Product security bad practices,'' Oct. 2024. [Online]. Available:
        \url{https://www.cisa.gov/resources-tools/resources/product-security-bad-practices}

        \bibitem{cwetop25}
        MITRE Corporation, ``CWE Top 25 most dangerous software weaknesses,'' 2024. [Online].
        Available: \url{https://cwe.mitre.org/top25/}

        \bibitem{dietz2012}
        W. Dietz, P. Li, J. Regehr, and V. Adve, ``Understanding integer overflow in C/C++,'' in
        \emph{Proc. 34th Int. Conf. Software Engineering (ICSE)}, Zurich, Switzerland, 2012,
        pp. 760--770.

        \bibitem{wang2012}
        X. Wang, H. Chen, Z. Jia, N. Zeldovich, and M. F. Kaashoek, ``Improving integer security
        for systems with KINT,'' in \emph{Proc. 10th USENIX Symp. Operating Systems Design and
        Implementation (OSDI)}, Hollywood, CA, USA, 2012, pp. 163--177.

        \bibitem{brumley2007}
        D. Brumley, T. Chiueh, R. Johnson, H. Lin, and D. Song, ``RICH: Automatically protecting
        against integer-based vulnerabilities,'' in \emph{Proc. Network and Distributed System
        Security Symp. (NDSS)}, San Diego, CA, USA, 2007.

        \bibitem{molnar2009}
        D. Molnar, X. C. Li, and D. A. Wagner, ``Dynamic test generation to find integer bugs in
        x86 binary Linux programs,'' in \emph{Proc. 18th USENIX Security Symp.}, Montreal,
        Canada, 2009, pp. 67--82.

        \bibitem{coker2013}
        Z. Coker and M. Hafiz, ``Program transformations to fix C integers,'' in \emph{Proc. 35th
        Int. Conf. Software Engineering (ICSE)}, San Francisco, CA, USA, 2013, pp. 792--801.

        \bibitem{zhang2010}
        C. Zhang, T. Wang, T. Wei, Y. Chen, and W. Zou, ``IntPatch: Automatically fix
        integer-overflow-to-buffer-overflow vulnerability at compile-time,'' in \emph{Proc. 15th
        European Symp. Research in Computer Security (ESORICS)}, Athens, Greece, 2010,
        pp. 71--86.

        \bibitem{ubsan}
        LLVM Project, ``UndefinedBehaviorSanitizer,'' Clang documentation. [Online]. Available:
        \url{https://clang.llvm.org/docs/UndefinedBehaviorSanitizer.html}

        \bibitem{gccmanual}
        Free Software Foundation, \emph{Using the GNU Compiler Collection (GCC)}, sec.\ ``Built-in
        functions to perform arithmetic with overflow checking'' and option \code{-ftrapv}.
        [Online]. Available: \url{https://gcc.gnu.org/onlinedocs/gcc/}

        \bibitem{ramey2019}
        R. Ramey, ``Boost.SafeNumerics,'' Boost C++ Libraries, 2019. [Online]. Available:
        \url{https://www.boost.org/doc/libs/release/libs/safe_numerics/}

        \bibitem{safeint}
        D. LeBlanc, ``SafeInt,'' Microsoft, commit 72f0745. [Online]. Available:
        \url{https://github.com/dcleblanc/SafeInt}

        \bibitem{intelsafe}
        Intel Corporation, ``safe arithmetic: A general purpose C++20 library for writing safe and
        bug-free code,'' pre-release, commit c3b6e94. [Online]. Available:
        \url{https://github.com/intel/safe-arithmetic}

        \bibitem{bounded}
        D. Stone, ``bounded::integer.'' [Online]. Available:
        \url{https://github.com/davidstone/bounded-integer}

        \bibitem{typesafe}
        J. M\"{u}ller, ``type\_safe.'' [Online]. Available:
        \url{https://github.com/foonathan/type_safe}

        \bibitem{rustnum}
        The Rust Project, ``Module std::num and the primitive integer types,'' \emph{The Rust
        Standard Library}. [Online]. Available: \url{https://doc.rust-lang.org/std/num/}

        \bibitem{spark}
        AdaCore, \emph{SPARK 2014 User's Guide}. [Online]. Available:
        \url{https://docs.adacore.com/spark2014-docs/html/ug/}

        \bibitem{why3}
        J.-C. Filli\^{a}tre and A. Paskevich, ``Why3: Where programs meet provers,'' in
        \emph{Proc. 22nd European Symp. Programming (ESOP)}, ser. LNCS, vol. 7792. Rome, Italy:
        Springer, 2013, pp. 125--128.

        \bibitem{rieuhelft2017}
        R. Rieu-Helft, C. March\'{e}, and G. Melquiond, ``How to get an efficient yet verified
        arbitrary-precision integer library,'' in \emph{Verified Software: Theories, Tools, and
        Experiments (VSTTE)}, ser. LNCS, vol. 10712. Springer, 2017, pp. 84--101.

        \bibitem{melquiond2018}
        G. Melquiond and R. Rieu-Helft, ``A Why3 framework for reflection proofs and its
        application to GMP's algorithms,'' in \emph{Proc. 9th Int. Joint Conf. Automated
        Reasoning (IJCAR)}, ser. LNCS, vol. 10900. Oxford, UK: Springer, 2018, pp. 178--193.

        \bibitem{glesner2006}
        S. Glesner, ``Finite integer computations: An algebraic foundation for their
        correctness,'' \emph{Formal Aspects of Computing}, vol. 18, no. 2, pp. 244--262, 2006.

        \bibitem{wilson1927}
        E. B. Wilson, ``Probable inference, the law of succession, and statistical inference,''
        \emph{J. Amer. Statist. Assoc.}, vol. 22, no. 158, pp. 209--212, 1927.

        \bibitem{mcnemar1947}
        Q. McNemar, ``Note on the sampling error of the difference between correlated proportions
        or percentages,'' \emph{Psychometrika}, vol. 12, no. 2, pp. 153--157, 1947.

        \bibitem{gbench}
        Google, ``Benchmark: A microbenchmark support library,'' v1.9.1. [Online]. Available:
        \url{https://github.com/google/benchmark}

        \bibitem{nvdapi}
        National Institute of Standards and Technology, ``National Vulnerability Database:
        Vulnerabilities API 2.0.'' [Online]. Available:
        \url{https://nvd.nist.gov/developers/vulnerabilities}

        \bibitem{p2900}
        J. Berne, T. Doumler, A. Krzemie\'{n}ski \emph{et al.}, ``Contracts for C++,'' ISO/IEC
        JTC1/SC22/WG21, P2900R14, 2025. [Online]. Available: \url{https://wg21.link/p2900r14}

        \bibitem{p3471}
        K. Varlamov and L. Dionne, ``Standard library hardening,'' ISO/IEC JTC1/SC22/WG21,
        P3471R1, 2024. [Online]. Available: \url{https://wg21.link/p3471r1}

        \bibitem{p2993}
        L. Valenty, ``Constrained numbers,'' ISO/IEC JTC1/SC22/WG21, P2993. [Online]. Available:
        \url{https://wg21.link/p2993}

        \bibitem{vazou2014}
        N. Vazou, E. L. Seidel, R. Jhala, D. Vytiniotis, and S. Peyton Jones, ``Refinement types
        for Haskell,'' in \emph{Proc. 19th ACM SIGPLAN Int. Conf. Functional Programming (ICFP)},
        Gothenburg, Sweden, 2014, pp. 269--282.

        \bibitem{safenumbers}
        M. Borland, ``SafeNumbers: Safer fundamental numeric types in C++20,'' The C++ Alliance,
        2026.
        [Online]. Available: \url{https://github.com/cppalliance/safe_numbers}, documentation at
        \url{https://develop.safe-numbers.cpp.al/}

    \end{thebibliography}

\end{document}
